\subsection{Frozen Versus Nonequilibrium Study\label{ch4:frozenVsNonequilibrium}}
This section investigates \gls{AO} properties at various Mach numbers for a
hypersonic vehicle and compares the results between frozen flow and
nonequilibrium flow conditions. \Cref{tab:frozenVsNonequilibriumStudy} summarizes
the key characteristics of each case analyzed. The gas composition if the same
composition from \Cref{tab:speciesStudyMassFraction}.
\begin{table}[H]
    \begin{center}
        \begin{tabular}{|c|c|c|c|c|c|c|}
            \hline
            Case & $M\Units{\;}$ & $T_{wall}\Units{K}$ &
            $T_{\infty}\Units{K}$ & $P_{\infty}\Units{Pa}$ &
            $h_0\Units{MJ}$ & $q\Units{kPa}$ \\
            \hline\hline
            $1R$\cite{Mackey2019} & $11.2$ & $1000$ & $293$ & $1200$ & $7.38$  & $105$ \\\hline
            $2R$\cite{Mackey2019} & $13.2$ & $1000$ & $293$ & $1200$ & $10.25$ & $146$ \\\hline
            $3R$\cite{Mackey2019} & $15.2$ & $1000$ & $293$ & $1200$ & $13.60$ & $193$ \\\hline
        \end{tabular}
    \caption{Chemistry Study \label{tab:frozenVsNonequilibriumStudy}}
    \end{center}
\end{table}

For this investigation the IRV-2 vehicle, shown in \Cref{fig:IRV2} is used.
The IRV-2 vehicle is a sphere-biconic with a nose radius of $1.95\Units{cm}$ and a total
length of $1.386 \Units{m}$. Six drag flaps are located at the aft end of
the vehicle and can be deployed to a maximum angle of $45$ degrees
\cite{Kuntz2001}. For this work, only axisymmetric calculations on the
IRV-2 nosetip are performed. \Cref{fig:IRV2Mesh} shows
the mesh used in this work. 
\begin{figure}[H]
    \centering
    \subfigure[IRV-2 vehicle]
    {
        \label{fig:IRV2}
        \includegraphics[height=0.5\textwidth]{chapter4/frozenVsNonequilibrium/irv2}
    }
    \subfigure[IRV-2 mesh]
    {
        \label{fig:IRV2Mesh}
        \includegraphics[height=0.5\textwidth]{chapter4/frozenVsNonequilibrium/1R_frozenMesh}
    }
    \caption{\label{fig:IRV2Images}IRV-2}
\end{figure}

\Cref{fig:frozenConvergence,fig:nonequilibriumConvergence} show the convergence
plots for cases described in \Cref{tab:frozenVsNonequilibrium}.
\begin{figure}[H]
    \centering
        \subfigure[Case $1R$]{
            \includegraphics[width=0.45\textwidth]{chapter4/frozenVsNonequilibrium/1R_frozen_convergence.pdf}}
        \subfigure[Case $2R$]{
            \includegraphics[width=0.45\textwidth]{chapter4/frozenVsNonequilibrium/2R_frozen_convergence.pdf}}
        \subfigure[Case $3R$]{
            \includegraphics[width=0.45\textwidth]{chapter4/frozenVsNonequilibrium/3R_frozen_convergence.pdf}}
        \caption{Convergence plot for frozen
        cases\label{fig:frozenConvergence}}
\end{figure}

\begin{figure}[H]
    \centering
        \subfigure[Case $1R$]{
            \includegraphics[width=0.45\textwidth]{chapter4/frozenVsNonequilibrium/1R_nonequilibrium_convergence.pdf}}
        \subfigure[Case $2R$]{
            \includegraphics[width=0.45\textwidth]{chapter4/frozenVsNonequilibrium/2R_nonequilibrium_convergence.pdf}}
        \subfigure[Case $3R$]{
            \includegraphics[width=0.45\textwidth]{chapter4/frozenVsNonequilibrium/3R_nonequilibrium_convergence.pdf}}
        \caption{Convergence plot for nonequilibrium
        cases\label{fig:nonequilibriumConvergence}}
\end{figure}


\subsubsection{Flow results}
\Cref{fig:contourTranslationTemperatureFrozen} presents contour plots of
translational temperature, while \Cref{fig:contourVibrationalTemperatureFrozen}
displays vibrational temperature contours at different Mach numbers for a
frozen flow. As expected, the post-shock temperature increases with Mach number
due to the rapid compression of the flow. In the case of a frozen flow, this
compression occurs as an adiabatic process, where kinetic energy is converted
into thermal energy, leading to a sudden temperature rise. Furthermore, higher
Mach numbers result in proportionally greater increases in both translational
and vibrational temperatures.
% Translational Temperatures Contours Frozen %
\begin{figure}[H]
    \centering
    \subfigure[Case
    $1R$]{\label{fig:frozenTr1R}\includegraphics[width=0.45\textwidth]{chapter4/frozenVsNonequilibrium/1R_frozenContourTtr.png}}
    \subfigure[Case
    $2R$]{\label{fig:frozenTr2R}\includegraphics[width=0.45\textwidth]{chapter4/frozenVsNonequilibrium/2R_frozenContourTtr.png}}
    \subfigure[Case
    $3R$]{\label{fig:frozenTr3R}\includegraphics[width=0.45\textwidth]{chapter4/frozenVsNonequilibrium/3R_frozenContourTtr.png}}
    \caption{Contour translational temperature for
    frozen cases from \Cref{tab:frozenVsNonequilibriumStudy}}
    \label{fig:contourTranslationTemperatureFrozen}
\end{figure}
% Translational Temperatures Contours Frozen %

% Vibrational Temperatures Contours Frozen %
\begin{figure}[H]
    \centering
    \subfigure[Case
    $1R$]{\label{fig:frozenTvr1R}\includegraphics[width=0.45\textwidth]{chapter4/frozenVsNonequilibrium/1R_frozenContourTvr.png}}
    \subfigure[Case
    $2R$]{\label{fig:frozenTvr2R}\includegraphics[width=0.45\textwidth]{chapter4/frozenVsNonequilibrium/2R_frozenContourTvr.png}}
    \subfigure[Case
    $3R$]{\label{fig:frozenTvr3R}\includegraphics[width=0.45\textwidth]{chapter4/frozenVsNonequilibrium/3R_frozenContourTvr.png}}
\caption{Contour vibrational temperature for
    frozen cases from
    \Cref{tab:frozenVsNonequilibriumStudy}\label{fig:contourVibrationalTemperatureFrozen}}
\end{figure}
% Vibrational Temperatures Contours Frozen %

\Cref{fig:contourTranslationTemperatureNonequilibrium} presents contour plots
of translational temperature, while
\Cref{fig:contourVibrationalTemperatureNonequilibrium} displays vibrational
temperature contours at different Mach numbers for a nonequilibrium flow. As
expected, the translational temperature increases with Mach number, primarily
at the leading edge of the geometry, where most of the kinetic energy of the
flow is converted into internal energy. However, the increase in vibrational
temperature is relatively lower due to nonequilibrium effects, as the process
is no longer adiabatic. Additionally, as expected, the frozen chemistry case
exhibits higher translational and vibrational temperatures since chemical
reactions do not remove energy from the flow.
% Translational Temperatures Contours Nonequilibrium %
\begin{figure}[H]
    \centering
    \subfigure[Case
    $1R$]{\label{fig:nonequilibriumTr1R}\includegraphics[width=0.45\textwidth]{chapter4/frozenVsNonequilibrium/1R_nonequilibriumContourTtr.png}}
    \subfigure[Case
    $2R$]{\label{fig:nonequilibriumTr2R}\includegraphics[width=0.45\textwidth]{chapter4/frozenVsNonequilibrium/2R_nonequilibriumContourTtr.png}}
    \subfigure[Case
    $3R$]{\label{fig:nonequilibriumTr3R}\includegraphics[width=0.45\textwidth]{chapter4/frozenVsNonequilibrium/3R_nonequilibriumContourTtr.png}}
\caption{Contour translational temperature for
            nonequilibrium cases from
            \Cref{tab:frozenVsNonequilibriumStudy}\label{fig:contourTranslationTemperatureNonequilibrium}}
\end{figure}
% Translational Temperatures Contours Nonequilibrium %
    

% Vibrational Temperatures Contours Nonequilibrium %
\begin{figure}[H]
    \centering
    \subfigure[Case
    $1R$]{\label{fig:nonequilibriumTvr1R}\includegraphics[width=0.45\textwidth]{chapter4/frozenVsNonequilibrium/1R_nonequilibriumContourTvr.png}}
    \subfigure[Case
    $2R$]{\label{fig:nonequilibriumTvr2R}\includegraphics[width=0.45\textwidth]{chapter4/frozenVsNonequilibrium/2R_nonequilibriumContourTvr.png}}
    \subfigure[Case
    $3R$]{\label{fig:nonequilibriumTvr3R}\includegraphics[width=0.45\textwidth]{chapter4/frozenVsNonequilibrium/3R_nonequilibriumContourTvr.png}}
\caption{Contour vibrational temperature for
            nonequilibrium cases from
            \Cref{tab:frozenVsNonequilibriumStudy}\label{fig:contourVibrationalTemperatureNonequilibrium}}
\end{figure}
% Vibrational Temperatures Contours Nonequilibrium %

\subsubsection{Optics results}
\Cref{fig:contourGDNonequilibriumR} shows the contour of \gls{GD} for a
nonequilibrium flow at three different Mach numbers. As expected, the highest
Mach number (\Cref{fig:nonequilibriumGD3R}) produces the highest \gls{GD}, since
post-shock flow properties, such as temperature, increase with Mach number.
Higher temperatures promote molecular dissociation, which significantly impacts
\gls{GD}. The highest \gls{GD} appears at the stagnation point, which is
reasonable because this region experiences the highest temperatures, leading to
the most intense dissociation effects.
% Gladstone Nonequilibrium % 
\begin{figure}[H]
    \centering
    \subfigure[Case
    $1R$]{\label{fig:nonequilibriumGD1R}\includegraphics[width=0.45\textwidth]{chapter4/frozenVsNonequilibrium/1R_nonequilibrium_gladstone_dale.png}}
    \subfigure[Case $2R$]{\label{fig:nonequilibriumGD2R}\includegraphics[width=0.45\textwidth]{chapter4/frozenVsNonequilibrium/2R_nonequilibrium_gladstone_dale.png}}
    \subfigure[Case $3R$]{\label{fig:nonequilibriumGD3R}\includegraphics[width=0.45\textwidth]{chapter4/frozenVsNonequilibrium/3R_nonequilibrium_gladstone_dale.png}}
    \caption{Contour \gls{GD} for nonequilibrium cases from
        \Cref{tab:frozenVsNonequilibriumStudy}}
    \label{fig:contourGDNonequilibriumR}
\end{figure}
% Gladstone Nonequilibrium % 

\Cref{fig:contourDiluteIndexOfRefractionNonequilibriumR} shows the contour
plots for the dilute index of refraction at different Mach numbers in a
nonequilibrium flow, while
\Cref{fig:contourDenseIndexOfRefractionNonequilibriumR} presents the contour
plot for the dense index of refraction under the same conditions. A noticeable
difference between the dense and dilute formulations of the index of refraction
is not observed. This suggests that the dilute formulation serves as a good
approximation, which is reasonable since air, even at high Mach numbers, is not
dense enough to produce significant deviations between the two formulations. As
expected, the highest variation in the index of refraction occurs along the
stagnation line, where temperature and density gradients are most pronounced
due to shock and compression effects in the nonequilibrium regime.
% Dilute Index of refraction Contours Nonequilibrium %
\begin{figure}[H]
    \centering
    \subfigure[Case $1R$]{\label{fig:nonequilibriumDiluteIndex1R}
    \includegraphics[width=0.45\textwidth]{chapter4/frozenVsNonequilibrium/1R_nonequilibrium_dilute_index.png}}
    \subfigure[Case $2R$]{\label{fig:nonequilibriumDiluteIndex2R}
    \includegraphics[width=0.45\textwidth]{chapter4/frozenVsNonequilibrium/2R_nonequilibrium_dilute_index.png}}
    \subfigure[Case $3R$]{\label{fig:nonequilibriumDiluteIndex3R}
    \includegraphics[width=0.45\textwidth]{chapter4/frozenVsNonequilibrium/3R_nonequilibrium_dilute_index.png}}
\caption{Contour dilute index of refraction for
            nonequilibrium cases from
            \Cref{tab:frozenVsNonequilibriumStudy}\label{fig:contourDiluteIndexOfRefractionNonequilibriumR}}
\end{figure}
% Dilute Index of refraction Contours Nonequilibrium %

% Dense Index of refraction Contours Nonequilibrium %
\begin{figure}[H]
    \centering
        \subfigure[Case
        $1R$]{\label{fig:nonequilibriumDenseIndex1R}\includegraphics[width=0.45\textwidth]{chapter4/frozenVsNonequilibrium/1R_nonequilibrium_dense_index.png}}
        \subfigure[Case
        $2R$]{\label{fig:nonequilibriumDenseIndex2R}\includegraphics[width=0.45\textwidth]{chapter4/frozenVsNonequilibrium/2R_nonequilibrium_dense_index.png}}
        \subfigure[Case
        $3R$]{\label{fig:nonequilibriumDenseIndex3R}\includegraphics[width=0.45\textwidth]{chapter4/frozenVsNonequilibrium/3R_nonequilibrium_dense_index.png}}
\caption{Contour dense index of refraction for
            nonequilibrium cases from
            \Cref{tab:frozenVsNonequilibriumStudy}\label{fig:contourDenseIndexOfRefractionNonequilibriumR}}
\end{figure}
% Dense Index of refraction Contours Nonequilibrium %

\Cref{fig:contourDiluteIndexOfRefractionFrozenR} shows the contour plots for the
dilute index of refraction at different Mach numbers in a frozen flow, while
\Cref{fig:contourDenseIndexOfRefractionFrozenR} presents the contour plot for
the dense index of refraction under the same conditions. A noticeable
difference between the dense and dilute formulations of the index of refraction
is not observed. This suggests that the dilute formulation serves as a good
approximation, which is reasonable since air, even at high Mach numbers, is not
dense enough to produce significant deviations between the two formulations.\\

As expected, the highest variation in the index of refraction occurs along the
stagnation line, where temperature and density gradients are most pronounced
due to shock and compression effects. However, a difference between the frozen
and nonequilibrium cases is observed: the frozen case appears to underpredict
the index of refraction. This is expected because, in the frozen case, the
density is underpredicted, and since the index of refraction is closely related
to the density through the \gls{GD}, a lower density results in a lower index
of refraction.
% Dilute Index of refraction Contours Frozen %
\begin{figure}[H]
    \centering
        \subfigure[Case $1R$]{
            \label{fig:frozenDiluteIndex1R}
            \includegraphics[width=0.45\textwidth]{chapter4/frozenVsNonequilibrium/1R_frozen_dilute_index.png}}
        \subfigure[Case $2R$]{
            \label{fig:frozenDiluteIndex2R}
            \includegraphics[width=0.45\textwidth]{chapter4/frozenVsNonequilibrium/2R_frozen_dilute_index.png}}
        \subfigure[Case $3R$]{
            \label{fig:frozenDiluteIndex3R}
            \includegraphics[width=0.45\textwidth]{chapter4/frozenVsNonequilibrium/3R_frozen_dilute_index.png}}
\caption{Contour dilute index of refraction for
            frozen cases from
            \Cref{tab:frozenVsNonequilibriumStudy}\label{fig:contourDiluteIndexOfRefractionFrozenR}}
\end{figure}

\begin{figure}[H]
    \centering
        \subfigure[Case $1R$]{
            \label{fig:frozenDenseIndex1R}
            \includegraphics[width=0.45\textwidth]{chapter4/frozenVsNonequilibrium/1R_frozen_dense_index.png}}
        \subfigure[Case $2R$]{
            \label{fig:frozenDenseIndex2R}
            \includegraphics[width=0.45\textwidth]{chapter4/frozenVsNonequilibrium/2R_frozen_dense_index.png}}
        \subfigure[Case $3R$]{
            \label{fig:frozenDenseIndex3R}
            \includegraphics[width=0.45\textwidth]{chapter4/frozenVsNonequilibrium/3R_frozen_dense_index.png}}
\caption{Contour dense index of refraction for
            frozen cases from
            \Cref{tab:frozenVsNonequilibriumStudy}\label{fig:contourDenseIndexOfRefractionFrozenR}}
\end{figure}
% Dense Index of refraction Contours Frozen %

\Cref{fig:contourDiluteElectricSusceptibilityNonequilibriumR} presents the contour
plots for the electric susceptibility at different Mach numbers for a
nonequilibrium flow. As expected, the electric susceptibility increases with
Mach number due to the higher post-shock temperatures and increased
dissociation effects that enhance the medium's polarization response leading to
an increase in electric susceptibility. Additionally, the
most significant variations in electric susceptibility occur along the
stagnation line, where temperature and density gradients are most pronounced.
This region experiences the highest degree of molecular excitation and
dissociation, further contributing to changes in the medium's electromagnetic
properties.
% Dilute Nonequilibrium Electric Susceptibility % 
\begin{figure}[H]
    \centering
        \subfigure[Case $1R$]{
            \label{fig:frozenDiluteIndex1R}
            \includegraphics[width=0.45\textwidth]{chapter4/frozenVsNonequilibrium/1R_nonequilibrium_susceptibility_dilute.png}}
        \subfigure[Case $2R$]{
            \label{fig:frozenDiluteIndex2R}
            \includegraphics[width=0.45\textwidth]{chapter4/frozenVsNonequilibrium/2R_nonequilibrium_susceptibility_dilute.png}}
        \subfigure[Case $3R$]{
            \label{fig:frozenDiluteIndex3R}
            \includegraphics[width=0.45\textwidth]{chapter4/frozenVsNonequilibrium/3R_nonequilibrium_susceptibility_dilute.png}}
\caption{Contour electric susceptibility for
            nonequilibrium cases from
            \Cref{tab:frozenVsNonequilibriumStudy}\label{fig:contourDiluteElectricSusceptibilityNonequilibriumR}}
\end{figure}
% Dilute Nonequilibrium Electric Susceptibility % 

\Cref{fig:contourDiluteElectricSusceptibilityFrozenR} presents the contour plots
for the electric susceptibility at different Mach numbers for a frozen flow. As
expected, the electric susceptibility increases with Mach number due to higher
post-shock temperatures and increased dissociation effects, which enhance the
medium's electric susceptibility response. Additionally, the most significant
variations in electric susceptibility occur along the stagnation line, where
temperature and density gradients are most pronounced. This region experiences
the highest degree of molecular excitation and dissociation, further
influencing the medium's electromagnetic properties. However, a notable
difference exists between the nonequilibrium and frozen cases: the frozen case
appears to underpredict the electric susceptibility, leading to an
underestimation of the medium's electromagnetic properties.
% Dilute Frozen Electric Susceptibility % 
\begin{figure}[H]
    \centering
        \subfigure[Case $1R$]{
            \label{fig:frozenElectricSusceptibility1R}
            \includegraphics[width=0.45\textwidth]{chapter4/frozenVsNonequilibrium/1R_frozen_susceptibility_dilute.png}}
        \subfigure[Case $2R$]{
            \label{fig:frozenElectricSusceptibility2R}
            \includegraphics[width=0.45\textwidth]{chapter4/frozenVsNonequilibrium/2R_frozen_susceptibility_dilute.png}}
        \subfigure[Case $3R$]{
            \label{fig:frozenElectricSusceptibility3R}
            \includegraphics[width=0.45\textwidth]{chapter4/frozenVsNonequilibrium/3R_frozen_susceptibility_dilute.png}}
\caption{Contour electric susceptibility for
            frozen cases from
            \Cref{tab:frozenVsNonequilibriumStudy}\label{fig:contourDiluteElectricSusceptibilityFrozenR}}
\end{figure}
% Dilute Frozen Electric Susceptibility % 

\subsubsection{Stagnation results}
\Cref{fig:stagnationTemperatures} shows the temperatures along the stagnation
line for different Mach numbers in both frozen and nonequilibrium flow
conditions. As expected, the shock standoff distance in the frozen case is
observed to be larger: approximately $2.3\Units{mm}$ for the slowest and medium
Mach numbers (Cases $1R$ and $2$), and $2\Units{mm}$ for the largest Mach
number (Case $3R$). On the other hand, the nonequilibrium case shows a shock
standoff distance of approximately $2\Units{mm}$ for the slow Mach number (Case
$1R$), $1.7\Units{mm}$ for the medium Mach number (Case $2R$), and
$1.5\Units{mm}$ for the largest Mach number (Case $3R$).\\

This difference arises because when chemistry effects are not considered, the
post-shock density is under-predicted (increasing the density ratio), and the
post-shock temperature is over-predicted (decreasing the temperature ratio).
The shock standoff scales with the density ratio across the shock times the
radius of the vehicle. When chemistry effects are neglected, the density ratio
across the shock is greater, and as a result, the shock standoff distance is
larger compared to when chemistry effects are taken into account.
% Stagnation Temperatures % 
\begin{figure}[H]
    \centering
        \subfigure[Case $1R$]{
            \includegraphics[width=0.45\textwidth]{chapter4/frozenVsNonequilibrium/1R_temperatures.pdf}}
        \subfigure[Case $2R$]{
            \includegraphics[width=0.45\textwidth]{chapter4/frozenVsNonequilibrium/2R_temperatures.pdf}}
        \subfigure[Case $3R$]{
            \includegraphics[width=0.45\textwidth]{chapter4/frozenVsNonequilibrium/3R_temperatures.pdf}}
    \caption{Stagnation temperatures: solid lines represent the nonequilibrium
    cases, while dashed lines represent the frozen cases}
    \label{fig:stagnationTemperatures}
\end{figure}
% Stagnation Temperatures % 

\Cref{fig:stagnationPressures} shows the pressure along the stagnation line at
different mach numbers for a frozen and nonequilibrium flow.
% Stagnation pressures % 
\begin{figure}[H]
    \centering
        \subfigure[Case $1R$]{
            \includegraphics[width=0.45\textwidth]{chapter4/frozenVsNonequilibrium/1R_pressures.pdf}}
        \subfigure[Case $2R$]{
            \includegraphics[width=0.45\textwidth]{chapter4/frozenVsNonequilibrium/2R_pressures.pdf}}
        \subfigure[Case $3R$]{
            \includegraphics[width=0.45\textwidth]{chapter4/frozenVsNonequilibrium/3R_pressures.pdf}}
    \caption{Pressures at the stagnation line}
    \label{fig:stagnationPressures}
\end{figure}
% Stagnation Pressures % 

\Cref{fig:stagnationMassFraction} illustrates the mass fraction of species
along the stagnation line. As expected, in the frozen flow case, there is no
change in mass fraction because there is no energy exchange with the internal
energy modes. However, in the nonequilibrium flow case, energy is exchanged
between internal energy modes, resulting in changes in the mass fractions of
chemical species. \\

For the nonequilibrium case, diatomic nitrogen (\ce{N2}) and diatomic oxygen
(\ce{O2}) dissociate partially, forming nitric oxide (\ce{NO}), atomic oxygen
(\ce{O}), and atomic nitrogen (\ce{N}) after the shock wave. The extent of this
dissociation increases with Mach number due to the higher post-shock
temperatures at greater Mach numbers. Higher temperatures provide the energy
required to overcome the dissociation energy of molecular species, leading to
greater production of \ce{NO}, \ce{O}, and \ce{N}.
% Stagnation mass fraction % 
\begin{figure}[H]
    \centering
        \subfigure[Case $1R$]{
            \includegraphics[width=0.45\textwidth]{chapter4/frozenVsNonequilibrium/1R_massFraction.pdf}}
        \subfigure[Case $2R$]{
            \includegraphics[width=0.45\textwidth]{chapter4/frozenVsNonequilibrium/2R_massFraction.pdf}}
        \subfigure[Case $3R$]{
            \includegraphics[width=0.45\textwidth]{chapter4/frozenVsNonequilibrium/3R_massFraction.pdf}}
    \caption{Stagnation mass fractions: solid lines represent the nonequilibrium
    cases, while dashed lines represent the frozen
    cases\label{fig:stagnationMassFraction}}
\end{figure}
% Stagnation mass fraction % 

\Cref{fig:stagnationSpeciesGD} illustrates the species \gls{GD} along the
stagnation line. As expected, in the frozen case, there is no change in the
species \gls{GD} because the mass fractions remain constant
(\Cref{fig:stagnationMassFraction}) due to the lack of exchange within internal
energy modes. However, in the nonequilibrium flow case, energy is exchanged
between internal energy modes, leading to a change in the mass fractions and,
consequently, a change in the species \gls{GD}. The species \gls{GD} follows a
trend very similar to that of the mass fractions; hence, the results presented
here closely resemble those in \Cref{fig:stagnationMassFraction}.
% Stagnation species GD % 
\begin{figure}[H]
    \centering
        \subfigure[Case $1R$]{
            \includegraphics[width=0.45\textwidth]{chapter4/frozenVsNonequilibrium/1R_speciesGD.pdf}}
        \subfigure[Case $2R$]{
            \includegraphics[width=0.45\textwidth]{chapter4/frozenVsNonequilibrium/2R_speciesGD.pdf}}
        \subfigure[Case $3R$]{
            \includegraphics[width=0.45\textwidth]{chapter4/frozenVsNonequilibrium/3R_speciesGD.pdf}}
    \caption{Stagnation species \gls{GD}: solid lines represent the nonequilibrium
    cases, while dashed lines represent the frozen
    cases\label{fig:stagnationSpeciesGD}}
\end{figure}
% Stagnation species GD % 

\Cref{fig:stagnationTotalGD} illustrates the total \gls{GD}. As expected, in
the frozen case, the total \gls{GD} remains constant due to the absence of
energy exchange between internal energy modes. In contrast, the nonequilibrium
case exhibits variation in the total \gls{GD}. Specifically, as the Mach number
increases, the total \gls{GD} also increases. This behavior is attributed to
the energy exchange between internal energy modes, which causes changes in the
thermodynamic and chemical state of the flow. The dependence of the total
\gls{GD} on the Mach number highlights the sensitivity of nonequilibrium flows
to the dynamic and thermal state of the system.
% Stagnation total GD % 
\begin{figure}[H]
    \centering
        \subfigure[Case $1R$]{
            \includegraphics[width=0.45\textwidth]{chapter4/frozenVsNonequilibrium/1R_totalGD.pdf}}
        \subfigure[Case $2R$]{
            \includegraphics[width=0.45\textwidth]{chapter4/frozenVsNonequilibrium/2R_totalGD.pdf}}
        \subfigure[Case $3R$]{
            \includegraphics[width=0.45\textwidth]{chapter4/frozenVsNonequilibrium/3R_totalGD.pdf}}
    \caption{Total \gls{GD} at the stagnation line\label{fig:stagnationTotalGD}}
\end{figure}
% Stagnation total GD % 

\Cref{fig:stagnationIndexOfRefraction} shows the stagnation index of refraction
for both the frozen and nonequilibrium cases at different Mach numbers. As
expected, the frozen case underpredicts the index of refraction and
overpredicts the standoff distance. This overprediction of the standoff
distance could significantly impact the beam quality of a seeker, as the
increased standoff distance may introduce additional phase distortions and beam
defocusing. Such effects could reduce the accuracy and effectiveness of optical
systems operating in hypersonic flow conditions. The nonequilibrium case, by
contrast, accounts for energy exchange between internal energy modes, resulting
in a more accurate representation of the index of refraction and associated
flow properties.
% Index of Refraction% 
\begin{figure}[H]
    \centering
        \subfigure[Case $1R$]{
            \includegraphics[width=0.45\textwidth]{chapter4/frozenVsNonequilibrium/1R_indexOfRefraction.pdf}}
        \subfigure[Case $2R$]{
            \includegraphics[width=0.45\textwidth]{chapter4/frozenVsNonequilibrium/2R_indexOfRefraction.pdf}}
        \subfigure[Case $3R$]{
            \includegraphics[width=0.45\textwidth]{chapter4/frozenVsNonequilibrium/3R_indexOfRefraction.pdf}}
    \caption{Index of refraction at the stagnation line\label{fig:stagnationIndexOfRefraction}}
\end{figure}
% Index of Refraction% 

\Cref{fig:stagnationElectricSusceptibility} compares the nonequilibrium and frozen
cases for electric susceptibility $\chi_e$ in a medium, as defined in
\Cref{eq:electricSusceptibility}. The behavior of electric susceptibility
follows a similar trend to the refractive index, leading the frozen case to
underpredict the medium’s response. A higher electric susceptibility indicates
that the material becomes more polarized when exposed to an external \gls{EM}
field. This increased polarization slows down signal propagation and enhances
signal attenuation. Consequently, the medium behaves more like an insulator,
impeding signal transmission and reducing power. In hypersonic flow conditions,
variations in electric susceptibility can significantly impact electromagnetic
wave-based systems, such as communication devices and radar seekers. Accurate
modeling of nonequilibrium effects is therefore essential for predicting and
mitigating signal degradation in high-speed aerospace applications.
% Electric Susceptibility % 
\begin{figure}[H]
    \centering
        \subfigure[Case $1R$]{
            \includegraphics[width=0.45\textwidth]{chapter4/frozenVsNonequilibrium/1R_electricSusceptibility.pdf}}
        \subfigure[Case $2R$]{
            \includegraphics[width=0.45\textwidth]{chapter4/frozenVsNonequilibrium/2R_electricSusceptibility.pdf}}
        \subfigure[Case $3R$]{
            \includegraphics[width=0.45\textwidth]{chapter4/frozenVsNonequilibrium/3R_electricSusceptibility.pdf}}
    \caption{Electric susceptibility at the stagnation
    line\label{fig:stagnationElectricSusceptibility}}
\end{figure}
% Electric Susceptibility % 
